\clearpage\section{System architecture}\label{sec:architecture}
This section describes the architecture as implemented in November 2026. It replaces the single-process design of the midterm report: the background worker now runs as its own service, chat streams its answers, uploads go to object storage, and the whole system is described as code for Amazon Web Services (AWS). Section~\ref{sec:results} gives the measurements behind each decision; Section~\ref{sec:progress} lists what remains unmeasured.
\subsection{Architectural style}
GENAILMS is a modular monolith. One backend codebase and one container image run as two services with different scaling signals: \code{api} serves requests and chat, and \code{worker} claims queued jobs for ingestion and module generation. Three supporting services sit beside them: \code{llm-gateway} (LiteLLM), \code{embedder} (Text Embeddings Inference with Qwen3) and, when the self-hosted option is chosen, \code{reranker} (Text Embeddings Inference with Jina). Microservices were considered and rejected: no part of the system is owned or deployed by a separate team, and splitting by business function would add network calls without adding an independent scaling signal. The split that matters is by load profile. Requests are short and numerous, generation jobs are long and wait mostly on the language model, and model serving is CPU or GPU bound.

Every service is stateless. Business data, vectors and the job queue live in PostgreSQL; files live in object storage; rate-limit counters live in Valkey, a Redis-compatible store, and may be lost. A task can therefore be added, replaced or killed without losing work, which is the precondition for horizontal scaling.
\subsection{Service responsibilities}
Table~\ref{tab:services} lists each service, the state it depends on and how it scales.
\begin{longtable}{|p{2.3cm}|p{5.4cm}|p{3.1cm}|p{3.0cm}|}\caption{Services, responsibilities and scaling signals.}\label{tab:services}\\\hline
Service & Responsibility & State & Scaling signal\\\hline\endfirsthead
\hline Service & Responsibility & State & Scaling signal\\\hline\endhead
\code{api} & Authentication, course and module operations, job creation, course chat over server-sent events (SSE), LTI launch & Runtime database role; no local state & Requests per target and CPU\\\hline
\code{worker} & Claims jobs with \code{FOR UPDATE SKIP LOCKED}; reads, chunks and embeds documents; generates modules; publishes to Moodle & Worker database role; temporary files only & Jobs ready in the queue\\\hline
\code{llm-gateway} & Holds provider keys, maps model aliases, retries and records served models & Configuration only & CPU (one task was enough in all tests)\\\hline
\code{embedder} & Qwen3-Embedding-0.6B vectors for chunks and queries; the model is baked into the image & Read-only model files & CPU\\\hline
\code{reranker} & Jina reranker v2 multilingual INT8; started only when Jina is the selected reranker & Read-only model files & Fixed count\\\hline
Frontend & React 19 and TypeScript single-page application, built to static files & None & Content delivery network\\\hline
\end{longtable}
The backend never sees provider keys: every model call, including embedding, goes through the gateway. The reranker is the exception to this rule, because both rerank options are called directly from the backend with a bounded number of concurrent requests per process.
\subsection{Local deployment}
Docker Compose runs the same services for development, with PostgreSQL and pgvector, Valkey, a Moodle 5.2 instance with the bridge plugin, and a Vite development server. On the measurement machine the embedder runs on its GPU. Compose builds the backend's \code{dev} image stage, which mounts the source and reloads on change; the production stage described below is what the load tests and AWS run. A separate load-test stack (\code{devops/loadtest}) runs the production image with its own database, a k6 load generator and a mock or real language model.
\subsection{Cloud deployment on AWS}\label{sec:aws}
Figure~\ref{fig:aws} shows the AWS deployment, written in Terraform under \code{infra/terraform}. It targets one region (Singapore, ap-southeast-1) and one environment per \code{terraform apply}, so the whole environment can be created before a demonstration or a measurement and destroyed afterwards. This is how the design fits the \$200 credit: nothing runs when nothing is being used.
\diagram{aws-deployment}{AWS deployment as written in Terraform. Dashed boxes are optional or outside AWS.}{aws}
\begin{longtable}{|p{3.2cm}|p{5.3cm}|p{5.4cm}|}\caption{AWS services, why each is used and what replaces it elsewhere.}\label{tab:aws-services}\\\hline
Need & Choice & Rejected alternative and lock-in\\\hline\endfirsthead
\hline Need & Choice & Rejected alternative and lock-in\\\hline\endhead
Run containers & ECS on Fargate: no servers to patch, one IAM role per task, scaling per service & EKS adds a control-plane fee and a Kubernetes ecosystem to operate; Lambda's 15-minute limit does not fit generation jobs. Lock-in: task definitions are rewritten, images are kept\\\hline
Database & RDS PostgreSQL 16 with pgvector; TLS required; backups and Multi-AZ are parameters & A separate vector database would need its own copy of course permissions, while every search is filtered by course. Lock-in: low, \code{pg_dump}\\\hline
Job queue & The existing PostgreSQL table; queue depth is published to CloudWatch & SQS would need an outbox, because job state must stay in the same transaction as the data it changes. At about 630 modules per hour the database used 2--9\% CPU\\\hline
Counters & ElastiCache Valkey, one node & Redis in a container saves little and is one more task to keep alive. Losing it only resets rate limits\\\hline
Files & S3 for uploads, deleted after embedding and expired after 7 days; S3 and CloudFront for the frontend & Shared file systems (EFS) cost more and keep a file-system model the tasks do not need\\\hline
Entry & CloudFront in front of an Application Load Balancer; the ALB accepts only requests carrying CloudFront's secret header & API Gateway's request time limit does not suit long chat streams\\\hline
Internal calls & ECS Service Connect, with the same names compose uses (\code{gateway}, \code{embedder}, \code{reranker}) & An internal ALB adds an hourly charge. No URL differs between local and cloud\\\hline
Secrets & SSM Parameter Store SecureString; provider keys reach only the gateway task & Secrets Manager charges per secret; automatic rotation is not yet needed\\\hline
Infrastructure & Terraform, validated in CI with \code{terraform test} against a mocked provider & CloudFormation and CDK tie the description to AWS\\\hline
\end{longtable}
Two cost decisions deviate from a textbook layout. First, there is no NAT gateway: tasks run in public subnets with public addresses but no inbound security-group rules, while RDS and Valkey sit in private subnets with no route to the Internet. A NAT gateway would cost about as much per hour as the whole demonstration compute. Second, the database is a single-AZ \code{db.t4g.medium}; Multi-AZ is one variable away for real operation. Terraform contains a check that fails the plan when the maximum task counts could open more database connections than the instance accepts (349 of 407 at the default limits).

The application code depends only on PostgreSQL, the Redis protocol, the S3 API and container images. What is specific to AWS lives in Terraform and the task definitions, so leaving AWS means rewriting the deployment layer, not the application.
\subsection{Images}
The backend Dockerfile has separate stages. The \code{prod} stage, which is the default build target and the one ECS runs, contains the installed dependencies, the application, the migrations and the three operator scripts (creating an administrator, registering an LTI platform and re-embedding chunks). It contains no package manager, no development dependencies, no tests and no reload watcher, and its files belong to root, so the process cannot modify its own code. Dependencies are compiled to bytecode at build time; this enlarges the image by about 160 MB but halved the time to import the application (3.7 s instead of 7.9 s), which every new task and worker process pays. All Python images use Python 3.14, the version continuous integration tests on.

The embedder image contains the five files of the pinned Qwen3 ONNX export (2.4 GB), each checked by SHA-256 at build time. A new embedder task therefore serves after about 8 s instead of first downloading the model from Hugging Face. Its vectors were compared with the running embedder's: cosine similarity 1.0, largest component difference $1.9\times10^{-7}$. This matters because stored chunk vectors are only comparable with vectors from the identical model export.
\subsection{Three load paths}
\subsubsection{Reads}
FastAPI opens the database session in one thread pool worker and runs the endpoint in another. Under 100 concurrent requests, every thread could end up waiting for a connection held by a request that was itself waiting for a thread; throughput fell from about 320 to 14 requests per second and did not recover. A concurrency-limit middleware now keeps the number of requests in service at or below the database pool size, which removed the stall.
\subsubsection{Module generation}
A module makes 22 to 23 model calls and about eight retrievals, and three quarters of its time used to be spent generating questions. Questions are now generated five at a time inside a module, retrieval for a question reranks 15 candidates instead of 30, and a busy reranker is retried with backoff instead of failing the job. Before this change, one failed rerank sent the whole job back to the queue, so model calls that had already succeeded were paid for again.

The external limits are the model provider's quota and, with the self-hosted option, the reranker. Rate-limit responses (HTTP 429) are handled per call: a background job waits and retries that one call, up to six times, and the gateway no longer takes the only deployment of a model out of service for 60 seconds after a single 429. This configuration is called L1 in Section~\ref{sec:results}. A shared limit on concurrent model calls across workers (L2) was designed but not adopted, because there is no measurement yet showing that it is needed. Requests from a waiting user are not retried this way: they fail quickly, so neither the user nor a database connection is held for minutes.
\subsubsection{Course chat}
Chat was the largest gap in the midterm design: a message held a request thread and a database connection while two model calls ran. The implemented path is:
\begin{enumerate}[nosep]
\item Save the question, commit and close the session.
\item Rewrite a follow-up question into a standalone one, using the last four messages, and search with the rewritten question. The answer step receives both the question as asked and the rewritten one, so a vague or unaccented follow-up is still understood; the answer language is read from the question as asked.
\item Search with at most eight concurrent searches per process, fewer than the database pool, and rerank with at most two attempts. If reranking fails, the hybrid order is kept; if the question cannot be embedded, the BM25 candidates are used alone. Generation does the opposite and fails the step, so the queue retries it with full retrieval.
\item Stream the answer without holding a thread, as SSE events \code{session}, \code{delta}, \code{done} or \code{error}, with a comment line every 15 s so that proxies do not close a silent connection. The ALB idle time-out (120 s) and CloudFront origin read time-out (60 s) are longer than this interval.
\item Read citations from markers such as \code{[S2]} in the streamed text, remove them from the displayed answer, and save the answer with its sources.
\end{enumerate}
A dedicated chat service was considered. It was not adopted, because the measurements show chat sharing \code{api} tasks without affecting other requests, and each extra service is one more thing to deploy and monitor. The limit \code{chat_max_streams} (100 per task by default) returns HTTP 503 before a task is overloaded.
\subsection{Reranking choice}\label{sec:rerank-choice}
Reranking on CPU was the shared bottleneck of both generation and chat (Section~\ref{sec:results}). Two options are implemented and selected by configuration. The self-hosted option is Jina reranker v2 INT8, whose licence permits only non-commercial use. The hosted option is Jev (\code{jev-1.14.0}), a rerank API from TypeSafe, called once per candidate chunk. Jev removed the bottleneck in every measurement and ranked at least as well. It also sends the query and the text of about 30 candidate chunks to a server in the United States, whose terms allow telemetry derived from the data and state no retention period. Jev is therefore implemented and measured but disabled by default until the university agrees to send course material there. A failed Jev call keeps the hybrid order, and each fallback is logged so that a CloudWatch alarm can report silent quality loss.
\subsection{State, failure and recovery}
Table~\ref{tab:state} lists where each kind of state lives and what happens when that place fails.
\begin{longtable}{|p{3.0cm}|p{3.6cm}|p{7.2cm}|}\caption{State and failure behaviour.}\label{tab:state}\\\hline
State or event & Where & Behaviour\\\hline\endfirsthead
\hline State or event & Where & Behaviour\\\hline\endhead
Business data, vectors & RDS & Point-in-time recovery from backups\\\hline
Jobs & PostgreSQL table & A worker refreshes a heartbeat once a minute, retrying a failed beat after 5 s; a sweeper requeues a job after five missed beats (5 minutes). An error from an unavailable dependency is retried after 30, 60, 120 and 240 s, so the five attempts outlast an outage of about seven minutes. A worker with a running job asks ECS for scale-in protection; when stopped, it gives the job 100 s to finish and then returns it to the queue without counting an attempt\\\hline
Sessions & Refresh tokens in PostgreSQL & Independent of tasks\\\hline
LTI \code{state} and \code{nonce} & PostgreSQL, 10-minute expiry & Not in Valkey: when Valkey is lost, rate limits fall back to per-task memory, but a login state cannot, because another task would not find it\\\hline
Rate limits & Valkey & Each task limits on its own until Valkey returns\\\hline
BM25 index & Task memory & A cache checked against a fingerprint on every search\\\hline
Readiness & \code{/ready} & Checks only the database, through a separate one-connection pool with a 2 s deadline. A database that accepted connections but did not answer once held \code{/ready} for 87 s; the deadline now bounds it\\\hline
Deployment & ECS & Migrations run as a one-off task of the new image before any service changes; a failed migration stops the deployment. A rollout whose tasks fail their health checks is rolled back by the ECS circuit breaker\\\hline
\end{longtable}
Running on RDS exposed one more problem before any deployment. PostgreSQL 16 lets only a role with \code{BYPASSRLS} change that attribute, and the RDS master user is not a superuser, so the first migration failed with \code{permission denied to alter role} when reproduced with a non-superuser owner. The worker is now exempted from row-level security by an explicit policy on each table instead of by a role attribute, in every environment, and a test fails when a new table with row-level security lacks that policy.
\subsection{Operations}
Logs are JSON lines in CloudWatch, kept for 14 days. The worker publishes \code{JobsReady}, \code{JobsRunning} and \code{OldestReadyJobAgeSeconds} once a minute in CloudWatch Embedded Metric Format, from a thread of its own so that it reports even when every job thread is busy. Because every worker process reports the same count, scaling and alarms read the maximum, not the sum. \code{api} scales on requests per target and CPU, \code{worker} on \code{JobsReady}, and the embedder on CPU. Alarms cover server errors and unhealthy targets at the ALB, a ready job waiting more than 15 minutes, database CPU, storage and connections, and Jev fallbacks or rejected Jev keys; AWS Budgets warns at \$50, \$100 and \$150. \code{infra/scripts/deploy.sh} builds and pushes images tagged by commit to repositories with immutable tags, runs the migration, applies Terraform and smoke-tests the public address.

The Terraform has passed \code{terraform validate} and plan tests against a mocked provider, but it has not yet been applied to a real AWS account. The first apply may still fail on account quotas or engine versions that only AWS reports.
